\chapter{L'apparato genitale maschile}

% ---------------------------------------------------------------------------
\section{Generalità}

A differenza di altre parti del corpo che funzionano in maniera continuativa, l'apparato genitale
è sempre presente ma resta ``inattivo'' fino alla pubertà (12--14 anni).

\rubrica{Costituzione}
\begin{itemize}
  \item \textbf{organi genitali} (gonadi): \textbf{testicoli};
  \item \textbf{vie genitali}: dai testicoli gli spermatozoi vengono veicolati all'esterno
    attraverso un sistema di condotti: tubuli retti, rete testis, condottini efferenti,
    \textbf{epididimo}, \textbf{condotto deferente}, \textbf{condotti eiaculatori}, \textbf{uretra}
    (in comune con l'apparato urinario);
  \item \textbf{grosse ghiandole annesse} alle vie genitali: \textbf{vescichette seminali},
    \textbf{prostata}, \textbf{ghiandole bulbo-uretrali}; versano il loro secreto nelle vie genitali
    durante l'eiaculazione;
  \item \textbf{genitali esterni}: \textbf{pene}, organo copulatore.
\end{itemize}

\rubrica{Funzione delle gonadi}
\begin{itemize}
  \item producono i \textbf{gameti} (spermatozoi/ovociti): cellule aploidi ($n$);
  \item producono \textbf{ormoni sessuali};
  \item i gameti maschili e femminili si possono fondere originando uno \textbf{zigote}, che darà
    origine a tutte le cellule di un nuovo individuo.
\end{itemize}

% ---------------------------------------------------------------------------
\section{Discesa dei testicoli}

Durante gli ultimi 2 mesi di gravidanza, normalmente i testicoli ``scendono'' dalla loro sede
originaria, la cavità addominale (ciascuno attraverso il canale inguinale del proprio lato),
raggiungendo lo scroto intorno al 7\textsuperscript{o} mese di gravidanza.

Un'arrestata discesa può avvenire a livello addominale o nel canale inguinale, per cause:
\begin{itemize}
  \item anatomiche;
  \item infiammatorie (orchite e periorchite fetale);
  \item genetiche (criptorchidismo familiare);
  \item inerenti al testicolo stesso (displasia o insufficiente risposta allo stimolo ormonale).
\end{itemize}

\textbf{Testicolo ritenuto}: posizionato al di fuori dello scroto, lungo il normale tragitto di
discesa. Può essere non visibile ma palpabile alla visita medica; dopo il riposizionamento manuale
nello scroto, tende a ritornare nella posizione iniziale appena cessata la manovra clinica.

\rubrica{Criptorchidismo}
Una delle più frequenti anomalie congenite dell'apparato urogenitale maschile.
\begin{itemize}
  \item \textbf{incidenza}: 3--5\% dei nati a termine, 9--30\% dei pre-termine;
  \item prima dei 5--6 mesi: non è necessario alcun trattamento particolare, in molti casi i
    testicoli scendono da soli nei primi mesi di vita;
  \item dopo i 6 mesi:
  \begin{enumerate}
    \item \textbf{terapia ormonale} a base di gonadotropine: efficace solo una volta su 4, non
      esente da effetti collaterali (possibili danni a lungo termine sulla spermatogenesi);
    \item \textbf{chirurgia precoce} (\textit{orchidopessi}): trattamento più seguito ed efficace,
      intervento mini invasivo in anestesia locale, piccolo taglietto scroto-inguinale per
      riportare il testicolo nella sede adeguata.
  \end{enumerate}
  \item un intervento tempestivo evita conseguenze in età adulta: riduzione della fertilità,
    insufficiente produzione di ormoni, fino alla degenerazione tumorale del testicolo.
\end{itemize}

% ---------------------------------------------------------------------------
\section{Ormoni sessuali}

Gli \textbf{ormoni sessuali} hanno funzioni fondamentali:
\begin{itemize}
  \item controllano lo sviluppo dei gameti e delle vie genitali;
  \item influiscono sullo sviluppo del corpo e di molti tessuti, sui caratteri sessuali secondari e
    sul comportamento.
\end{itemize}

Gli \textbf{ormoni sessuali maschili} (\textbf{androgeni}) sono prodotti soprattutto dal testicolo
e dalla corteccia surrenale.

La produzione di spermatozoi e la produzione di ormoni nei testicoli sono funzioni svolte da
popolazioni di cellule completamente \textbf{diverse}.

% ---------------------------------------------------------------------------
\section{Scroto e testicolo}

\textbf{Scroto}: sacca formata dalla cute e dalla fascia superficiale; si trova alla radice delle
cosce e alla base del pene. Un setto mediano lo divide in due cavità, una per ciascun testicolo.

\rubrica{Testicolo}
\begin{itemize}
  \item dimensioni: $\sim$4 cm di lunghezza $\times$ 2,5 cm di diametro;
  \item forma di ovoide: polo superiore, polo inferiore, margine posteriore e anteriore, faccia
    laterale e mediale;
  \item circondato da 2 tuniche: \textbf{tunica vaginale} e \textbf{tunica albuginea} (chiara,
    fibrosa);
  \item sospeso nello scroto: al polo superiore dal \textbf{funicolo} (cordone spermatico), al polo
    inferiore dal \textbf{legamento scrotale}.
\end{itemize}

La \textbf{tunica vaginale} è una membrana sierosa di derivazione peritoneale, organizzata in 2
foglietti; più esternamente vi è la \textbf{fascia spermatica}, un rivestimento fibroso del
testicolo.

Dalla tunica albuginea diparte un setto (\textbf{mediastino}, sul margine posteriore del
testicolo) che divide il testicolo in 250--300 compartimenti a forma di spicchi, detti
\textbf{lobuli}.
\begin{itemize}
  \item ciascun lobulo contiene da 1 a 4 \textbf{tubuli seminiferi} contorti, che producono i
    gameti (spermatozoi).
\end{itemize}

% ---------------------------------------------------------------------------
\section{Struttura del testicolo}

Il parenchima del testicolo è dato da: tubuli seminiferi contorti (4--5 per lobulo) e cellule
interstiziali.
\begin{itemize}
  \item i tubuli seminiferi contorti hanno pareti muscolo-elastiche e, nel lume, un epitelio
    stratificato detto \textbf{epitelio germinativo};
  \item nell'epitelio germinativo troviamo le cellule responsabili della gametogenesi (organizzate
    in più strati) e le \textbf{cellule del Sertoli}.
\end{itemize}

\rubrica{Funzioni delle cellule del Sertoli}
\begin{itemize}
  \item barriera ematica-testicolare;
  \item trasporto di nutrienti alle cellule germinative;
  \item trasporto delle cellule germinative dalla base al lume del tubulo;
  \item secrezione di fluido di trasporto;
  \item secrezione di messaggeri chimici regolatori della spermatogenesi.
\end{itemize}

% ---------------------------------------------------------------------------
\section{Vie spermatiche}

\begin{itemize}
  \item tubuli retti;
  \item rete testis;
  \item condottini efferenti;
  \item epididimo;
  \item dotto deferente;
  \item uretra.
\end{itemize}
Le vie spermatiche conducono lo sperma dalle gonadi all'esterno del corpo e permettono la
maturazione degli spermatozoi stessi; la loro funzionalità è regolata dagli ormoni sessuali
maschili.

I \textbf{tubuli seminiferi}, che producono gli spermatozoi, convergono a costituire un
\textbf{tubulo retto} per ogni lobulo.
\begin{itemize}
  \item il tubulo retto è un condotto rettilineo che conduce lo sperma alla \textbf{rete testis};
  \item la rete testis è una rete di tubuli da cui dipartono i \textbf{condottini efferenti};
  \item i condottini efferenti escono dal testicolo (attraverso il mediastino, margine posteriore)
    e conducono lo sperma nell'\textbf{epididimo}, localizzato sulla superficie esterna del
    testicolo.
\end{itemize}

% ---------------------------------------------------------------------------
\section{Pubertà e caratteri sessuali secondari}

Alla pubertà, il testicolo produce una quantità elevata di \textbf{testosterone}, che determina:
\begin{itemize}
  \item il \textbf{differenziamento puberale delle gonadi} e la ripresa della gametogenesi;
  \item l'\textbf{induzione dei caratteri sessuali secondari}:
  \begin{itemize}
    \item comparsa del pelo pubico;
    \item ingrandimento dei testicoli e del pene;
    \item crescita della barba;
    \item cambiamenti della voce;
    \item diversa distribuzione del grasso e dei muscoli.
  \end{itemize}
\end{itemize}

% ---------------------------------------------------------------------------
\section{Spermatogenesi}

La spermatogenesi avviene nei testicoli alla pubertà, alla ripresa dell'attività ormonale. È
suddivisa in 3 fasi:
\begin{enumerate}
  \item \textbf{proliferazione} $\rightarrow$ 16 giorni;
  \item \textbf{meiosi} $\rightarrow$ 24 giorni;
  \item \textbf{spermiogenesi} $\rightarrow$ 24 giorni.
\end{enumerate}
Totale: 64 giorni.

\subsection{I fase: proliferazione}
Compartimento basale (16 giorni). Nel testicolo esistono 2 popolazioni di spermatogoni:
\begin{itemize}
  \item \textbf{spermatogoni di tipo A} (scuri, cellule staminali): possono autorinnovarsi o
    riprendere la spermatogenesi, originando gli spermatogoni di tipo A1;
  \item \textbf{spermatogoni di tipo A1} (chiari): iniziano la fase proliferativa (I mitosi).
\end{itemize}

\subsection{II fase: meiosi}
Compartimento adluminale (24 giorni). Gli spermatociti I migrano nel compartimento adluminale e
compiono le due divisioni meiotiche. Gli spermatidi sono dotati di espressione genica aploide,
diversa da quella degli spermatociti I e II.

\subsection{III fase: spermiogenesi}
24 giorni. Rimodellamento cellulare che coinvolge nucleo e citoplasma e trasforma gli spermatidi
in spermatozoi (differenziamento degli spermatozoi). Si formano: \textbf{testa}, \textbf{acrosoma},
\textbf{coda}, \textbf{corpo residuo}.

% ---------------------------------------------------------------------------
\section{Struttura dello spermatozoo}

Spermatozoo lungo $\sim$60 $\mu$m; testa 4--5 $\mu$m.

\rubrica{Testa}
\begin{itemize}
  \item \textbf{nucleo}: contiene la cromatina;
  \item \textbf{acrosoma}: lisosoma primario che incappuccia il nucleo per 2/3 della sua
    lunghezza; compito: aprirsi un varco nella parete dell'ovulo.
\end{itemize}

\rubrica{Coda}
Costituita da un flagello molto lungo ($\sim$50--60 $\mu$m); presenta caratteristiche
morfologiche differenti lungo il suo decorso, in 4 parti:
\begin{enumerate}
  \item \textbf{collo}: diametro leggermente maggiore di qualsiasi altra parte della coda,
    contiene i residui citoplasmatici dello spermatide;
  \item \textbf{tratto intermedio}: costituito dalle 9 colonne segmentate che circondano il
    flagello interno, a struttura microtubulare tipica;
  \item \textbf{tratto principale}: caratterizzato dal flagello rivestito dalle colonne striate;
  \item \textbf{tratto terminale}: struttura microtubulare rivestita dalla membrana plasmatica.
\end{enumerate}

% ---------------------------------------------------------------------------
\section{Anomalie della spermatogenesi}

\begin{itemize}
  \item errori nella meiosi $\rightarrow$ gameti con un numero errato di cromosomi
    (\textbf{aneuploidie});
  \item variazioni nel numero o nella forma $\rightarrow$ \textbf{azospermia}, \textbf{oligospermia}.
\end{itemize}

\rubrica{Cause}
\begin{itemize}
  \item genetiche;
  \item infezioni croniche, prostatiti;
  \item criptorchidismo (sensibilità al testosterone);
  \item ipogonadismo endocrino;
  \item anticorpi antispermatozoo.
\end{itemize}

% ---------------------------------------------------------------------------
\section{Epididimo}

\textbf{Epididimo}: struttura a forma di virgola, lunghezza $\sim$4 cm, che appoggia sul polo
superiore del testicolo e sulla parete postero-laterale del testicolo stesso.
\begin{itemize}
  \item si distinguono \textbf{testa}, \textbf{corpo} e \textbf{coda};
  \item la testa (sul polo superiore del testicolo) riceve i condottini efferenti;
  \item la coda si continua verso l'alto nel condotto deferente;
  \item struttura: condotto avvolto e tortuoso che, disteso, misura $\sim$6 metri.
\end{itemize}

L'epididimo è rivestito da una tonaca mucosa con epitelio pseudostratificato. Alcune cellule
dell'epitelio hanno lunghi microvilli, non mobili, che assorbono l'eccesso di fluido testicolare
e rilasciano nutrienti allo sperma nel lume.

Gli spermatozoi, immaturi e quasi privi di motilità, lasciano i testicoli e percorrono lentamente
l'epididimo; durante questo percorso ($\sim$20 giorni) acquistano motilità.
\begin{itemize}
  \item l'epididimo ha una parete con muscolatura liscia: nell'eiaculazione essa si contrae e lo
    sperma è spinto dalla porzione terminale (coda) al successivo segmento di trasporto, il
    \textbf{dotto deferente};
  \item lo sperma può restare nell'epididimo per diversi mesi; se vi resta più a lungo, le cellule
    spermatiche vengono fagocitate dalle cellule di rivestimento.
\end{itemize}

% ---------------------------------------------------------------------------
\section{Dotto deferente}

\textbf{Dotto} (o vaso) \textbf{deferente}: lungo 40--45 cm; sale dall'epididimo, attraverso il
canale inguinale, fino alla cavità pelvica. In questo percorso è contenuto entro il
\textbf{cordone spermatico} (funicolo), a cui è sospeso il testicolo. Fa parte del cordone
spermatico ed è irrorato dall'\textbf{arteria deferenziale}.

\rubrica{Irrorazione sanguigna del testicolo}
\begin{itemize}
  \item i testicoli sono raggiunti dalla lunga \textbf{arteria testicolare}, ramo dell'aorta
    discendente addominale;
  \item il sangue refluo è raccolto da vasi venosi che si organizzano attorno all'arteria
    testicolare in una rete (\textbf{plesso pampiniforme}); il plesso sottrae calore dal sangue
    dell'arteria testicolare, che entra così nei testicoli raffreddato, aiutando a mantenerli alla
    temperatura ottimale.
\end{itemize}

% ---------------------------------------------------------------------------
\section{Vescicole seminali}

\textbf{Vescicole seminali}: appoggiano sulla parete posteriore della vescica; ghiandole lunghe
5--7 cm, di forma simile a un piccolo dito.

Il 60\% del volume del liquido seminale è secreto dalle vescicole seminali.
\begin{itemize}
  \item fluido alcalino, viscoso, giallastro, contenente fruttosio, acido ascorbico, un enzima
    coagulante e prostaglandine;
  \item da ciascuna vescicola seminale diparte un dotto che confluisce nel dotto deferente
    collaterale, a formare il \textbf{dotto eiaculatorio}.
\end{itemize}

% ---------------------------------------------------------------------------
\section{Prostata}

\textbf{Prostata}: ghiandola impari, di tessuto fibroso, posta al di sotto della vescica, attorno
all'uretra; racchiusa da una spessa capsula connettivale.
\begin{itemize}
  \item costituita da 20--30 corpi ghiandolari tubulo-alveolari, immersi in una massa di
    muscolatura liscia e connettivo;
  \item compito principale: produrre e immagazzinare il liquido seminale, rilasciato durante
    l'eiaculazione.
\end{itemize}

\rubrica{Morfologia}
\begin{itemize}
  \item forma piramidale, con base in alto e apice in basso, simile a una castagna; talvolta a
    mezzaluna o, in caso di ipertrofia, a ciambella;
  \item misure: $4 \times 3 \times 2$ cm; peso 10--20 g;
  \item palpabile mediante esame rettale, essendo collocata $\sim$5 cm anteriormente al retto e
    all'ano.
\end{itemize}

Iperplasia prostatica benigna e cancro alla prostata (la forma di tumore più frequente nel sesso
maschile) sono tra le malattie più frequenti.

Il \textbf{secreto prostatico} costituisce circa un terzo del volume del fluido seminale; nutre e
attiva gli spermatozoi. Si riversa nell'uretra, attraverso numerosi dotti, all'atto
dell'eiaculazione, quando la muscolatura liscia della prostata si contrae.

% ---------------------------------------------------------------------------
\section{Ghiandole accessorie}

\begin{itemize}
  \item \textbf{vescicole seminali} (pari);
  \item \textbf{ghiandole bulbo-uretrali} (pari);
  \item \textbf{prostata} (impari).
\end{itemize}
Producono la maggior parte del volume del liquido seminale.

% ---------------------------------------------------------------------------
\section{Pene}

\textbf{Pene}: organo predisposto all'accoppiamento, finalizzato a veicolare lo sperma nel tratto
riproduttivo femminile. Consta di \textbf{radice}, \textbf{corpo}, \textbf{glande}.
\begin{itemize}
  \item la radice è fissa, il corpo è libero e all'estremità si allarga nel glande;
  \item il glande è rivestito di cute, che si continua con la cute del corpo del pene; nel glande
    la cute origina una piega, il \textbf{prepuzio};
  \item internamente il pene è percorso in tutta la sua lunghezza dall'\textbf{uretra}; contiene
    inoltre tre corpi cilindrici, gli \textbf{organi erettili};
  \item gli organi erettili sono costituiti da una rete spugnosa di tessuto connettivo e
    muscolatura liscia, bucherellata e attraversata da vasi sanguigni.
\end{itemize}

% ---------------------------------------------------------------------------
\section{Il perineo}

Il pavimento della cavità pelvica è dato da strati muscolari: il \textbf{perineo}.
\begin{itemize}
  \item membrana muscolare a forma di losanga, tesa tra la sinfisi pubica (anteriormente), il
    coccige (posteriormente) e le tuberosità ischiatiche (lateralmente);
  \item muscoli del perineo: \textbf{ischiocavernoso}, \textbf{bulbospongioso}, \textbf{trasverso
    superficiale del perineo}.
\end{itemize}

% ---------------------------------------------------------------------------
\section{Uretra maschile}

\textbf{Uretra}: porzione terminale del sistema di trasporto dello sperma (comune all'apparato
urinario e all'apparato genitale).
\begin{itemize}
  \item origine: vescica, orifizio uretrale interno;
  \item termine: glande, orifizio uretrale esterno.
\end{itemize}

\rubrica{3 regioni dell'uretra}
\begin{enumerate}
  \item \textbf{uretra prostatica}: circondata dalla ghiandola prostatica; vi si aprono i dotti
    eiaculatori destro e sinistro;
  \item \textbf{uretra membranosa}: attraversa il diaframma urogenitale (parete muscolare striata);
  \item \textbf{uretra cavernosa}: circondata dai corpi cavernosi del pene; vi si aprono i condotti
    (destro e sinistro) delle ghiandole bulbo-uretrali.
\end{enumerate}

% ---------------------------------------------------------------------------
\section{Sintesi: organi e funzioni}

\begin{table}[htbp]
  \centering \small \renewcommand{\arraystretch}{1.3}
  \begin{tabularx}{\textwidth}{|l|X|}
    \hline
    \textbf{Organo} & \textbf{Funzione} \\
    \hline
    Testicolo & produzione di cellule germinali (gameti); produzione di ormoni \\
    \hline
    Epididimo & organo di raccolta degli spermatozoi \\
    \hline
    Dotto deferente & organo di trasporto degli spermatozoi \\
    \hline
    Uretra & organo di trasporto degli spermatozoi e organo urinario \\
    \hline
    Ghiandole genitali accessorie (prostata, vescichette seminali, ghiandole bulbo-uretrali)
      & produzione del liquido seminale \\
    \hline
    Pene & organo di copulazione e organo urinario \\
    \hline
  \end{tabularx}
\end{table}
